\chapter{2004年全国III卷（理）}
\section{单选题}
\begin{problem}
设集合 \(M = \{(x, y)\mid x^2 + y^2 = 1, x \in \R, y \in \R\}\)，\(N = \{(x, y)\mid x^2 - y = 0, x \in \R, y \in \R\}\)，则集合 \(M \cap N\) 中元素的个数为（\quad）

\choices
    {\(1\)}
    {\(2\)}
    {\(3\)}
    {\(4\)}
\end{problem}

\begin{problem}
函数 \(y = \left|\sin \frac{x}{2}\right|\) 的最小正周期是（\quad）

\choices
    {\(\frac{\pi}{2}\)}
    {\(\pi\)}
    {\(2\pi\)}
    {\(4\pi\)}
\end{problem}

\begin{problem}
设数列 \(\{a_{n}\}\) 是等差数列，且 \(a_{2} = -6\)，\(a_{8} = 6\)，\(S_{n}\) 是数列 \(\{a_{n}\}\) 的前 \(n\) 项和，则（\quad）

\choices
    {\(S_{4} < S_{5}\)}
    {\(S_{4} = S_{5}\)}
    {\(S_{6} < S_{5}\)}
    {\(S_{6} = S_{5}\)}
\end{problem}

\begin{problem}
圆 \(x^{2} + y^{2} - 4x = 0\) 在点 \(P(1, \sqrt{3})\) 处的切线方程为（\quad）

\choices
    {\(x + \sqrt{3} y - 2 = 0\)}
    {\(x + \sqrt{3} y - 4 = 0\)}
    {\(x - \sqrt{3} y + 4 = 0\)}
    {\(x - \sqrt{3} y + 2 = 0\)}
\end{problem}

\begin{problem}
函数 \(y = \sqrt{\log_{\frac{1}{2}}(x^2 - 1)}\) 的定义域为（\quad）

\choices
    {\(\left[-\sqrt{2}, -1\right) \cup (1, \sqrt{2}]\)}
    {\((-\sqrt{2}, -1) \cup (1, \sqrt{2})\)}
    {\(\left[-2, -1\right) \cup (1, 2]\)}
    {\((-2, -1) \cup (1, 2)\)}
\end{problem}

\begin{problem}
设复数 \(z\) 的辐角的主值为 \(\frac{2\pi}{3}\)，虚部为 \(\sqrt{3}\)，则 \(z^2 =\)（\quad）

\choices
    {\(-2 - 2\sqrt{3}\i\)}
    {\(-2\sqrt{3} - 2\i\)}
    {\(2 + 2\sqrt{3}\i\)}
    {\(2\sqrt{3} + 2\i\)}
\end{problem}

\begin{problem}
设双曲线的焦点在 \(x\) 轴上，两条渐近线为 \(y = \pm \frac{1}{2} x\)，则该双曲线的离心率 \(e =\)（\quad）

\choices
    {\(5\)}
    {\(\sqrt{5}\)}
    {\(\frac{\sqrt{5}}{2}\)}
    {\(\frac{5}{4}\)}
\end{problem}

\begin{problem}
不等式 \(1 < |x + 1| < 3\) 的解集为（\quad）

\choices
    {\((0,2)\)}
    {\((-2,0) \cup (2,4)\)}
    {\((-4,0)\)}
    {\((-4,-2) \cup (0,2)\)}
\end{problem}

\begin{problem}
正三棱锥的底面边长为 \(2\)，侧面均为直角三角形，则此三棱锥的体积为（\quad）

\choices
    {\(\frac{2\sqrt{2}}{3}\)}
    {\(\sqrt{2}\)}
    {\(\frac{\sqrt{2}}{3}\)}
    {\(\frac{4\sqrt{2}}{3}\)}
\end{problem}

\begin{problem}
在 \(\triangle ABC\) 中，\(AB = 3\)，\(BC = \sqrt{13}\)，\(AC = 4\)，则边 \(AC\) 上的高为（\quad）

\choices
    {\(\frac{3\sqrt{2}}{2}\)}
    {\(\frac{3\sqrt{3}}{2}\)}
    {\(\frac{3}{2}\)}
    {\(3\sqrt{3}\)}
\end{problem}

\begin{problem}
设函数 \(f(x) = \begin{cases}
(x + 1)^2, & x < 1,\\
4 - \sqrt{x - 1}, & x \geqslant 1,
\end{cases}\) 则使得 \(f(x) \geqslant 1\) 的自变量 \(x\) 的取值范围为（\quad）

\choices
    {\((-\infty, -2] \cup [0, 10]\)}
    {\((-\infty, -2] \cup [0, 1]\)}
    {\((-\infty, -2] \cup [1, 10]\)}
    {\([-2,0] \cup [1,10]\)}
\end{problem}

\begin{problem}
将 \(4\text{ 名}\) 教师分配到 \(3\text{ 所}\) 中学任教，每所中学至少 \(1\text{ 名}\)，则不同的分配方案共有（\quad）

\choices
    {\(12\text{ 种}\)}
    {\(24\text{ 种}\)}
    {\(36\text{ 种}\)}
    {\(48\text{ 种}\)}
\end{problem}

\section{填空题}
\begin{problem}
用平面 \(\alpha\) 截半径为 \(R\) 的球，如果球心到平面 \(\alpha\) 的距离为 \(\frac{R}{2}\)，那么截得小圆的面积与球的表面积的比值为\fillinblank{}.
\end{problem}

\begin{problem}
函数 \(y = \sin x + \sqrt{3}\cos x\) 在区间 \(\left[0, \frac{\pi}{2}\right]\) 上的最小值为\fillinblank{}.
\end{problem}

\begin{problem}
已知函数 \(y = f(x)\) 是奇函数，当 \(x \geqslant 0\) 时，\(f(x) = 3^x - 1\)，设 \(f(x)\) 的反函数是 \(y = g(x)\)，则 \(g(-8) =\)\fillinblank{}.
\end{problem}

\begin{problem}
设 \(P\) 是由曲线 \(y^{2} = 4(x - 1)\) 上的一个动点，则点 \(P\) 到点 \((0,1)\) 的距离与点 \(P\) 到 \(y\) 轴的距离之和的最小值为\fillinblank{}.
\end{problem}

\section{解答题}
\begin{problem}
已知 \(\alpha\) 为锐角，且 \(\tan \alpha = \frac{1}{2}\)，求 \(\frac{\sin 2\alpha \cos \alpha - \sin \alpha}{\sin 2\alpha \cos 2\alpha}\) 的值.
\end{problem}

\begin{problem}
解方程：\(4^{x} + |1 - 2^{x}| = 11\).
\end{problem}

\begin{problem}
某村计划建造一个室内面积为 \(800\mathrm{~m}^{2}\) 的矩形蔬菜温室. 在温室内，沿左，右两侧与后侧内墙各保留 \(1\mathrm{~m}\) 宽的通道，沿前侧内墙保留 \(3\mathrm{~m}\) 宽的空地. 当矩形温室的边长各为多少时？蔬菜的种植面积最大，最大种植面积是多少？
\end{problem}

\begin{problem}
三棱锥 \(P - ABC\) 中，侧面 \(PAC\) 与底面 \(ABC\) 垂直，\(PA = PB = PC = 3\).

\begin{enumerate}
\item 求证：\(AB \perp BC\)；
\item 设 \(AB = BC = 2\sqrt{3}\)，求 \(AC\) 与平面 \(PBC\) 所成角的大小.
\end{enumerate}

\bitmapfigure[max width=0.42\linewidth]{img/2004/national_paper_3_science/q20_fig1.png}
\end{problem}

\begin{problem}
设椭圆 \(\frac{x^2}{m+1} + y^2 = 1\) 的两个焦点是 \(F_{1}(-c,0)\) 与 \(F_{2}(c,0)\)（\(c > 0\)），且椭圆上存在一点 \(P\)，使得直线 \(PF_{1}\) 与 \(PF_{2}\) 垂直.

\begin{enumerate}
\item 求实数 \(m\) 的取值范围；
\item 设 \(L\) 是相应于焦点 \(F_{2}\) 的准线，直线 \(PF_{2}\) 与 \(L\) 相交于点 \(Q\). 若 \(\left|\frac{QF_{2}}{PF_{2}}\right| = 2 - \sqrt{3}\)，求直线 \(PF_{2}\) 的方程.
\end{enumerate}
\end{problem}

\begin{problem}
已知数列 \(\{a_{n}\}\) 的前 \(n\) 项和 \(S_{n}\) 满足 \(S_{n} = 2a_{n} + (-1)^{n}\)，\(n \geqslant 1\).

\begin{enumerate}
\item 写出数列 \(\{a_{n}\}\) 的前三项 \(a_{1}\)，\(a_{2}\)，\(a_{3}\)；
\item 求数列 \(\{a_{n}\}\) 的通项公式；
\item 证明：对任意的整数 \(m > 4\)，有 \(\frac{1}{a_4} + \frac{1}{a_5} + \cdots + \frac{1}{a_m} < \frac{7}{8}\).
\end{enumerate}
\end{problem}
